\documentclass[a4paper,twoside]{article}

\usepackage{morewrites}

\usepackage[svgnames,dvipsnames,x11names]{xcolor}
\usepackage{graphicx}
\usepackage[english]{babel}
\usepackage[colorlinks=true, linkcolor=NavyBlue, urlcolor=NavyBlue, citecolor=NavyBlue]{hyperref}
\usepackage[a4paper]{geometry}
\usepackage{ts-fonts}
\usepackage{amsmath}
\usepackage{float}
\usepackage{seqsplit}
\usepackage{ts-code}
\usepackage{ts-glossary}

\usepackage{lastpage}
\usepackage{refcount}
\makeatletter
\def\ps@plain{
  \let\@oddhead\@empty
  \let\@evenhead\@empty
  \def\@oddfoot{
    \hfil
    \ifnum\getpagerefnumber{LastPage}>1\relax
      \thepage
    \fi
    \hfil}
  \let\@evenfoot\@oddfoot
}
\makeatother

\title{Abbreviations / Acronyms}
\author{}\date{}

\begin{document}

\maketitle

\thispagestyle{plain}
Abbreviations are a lightweight mechanism to define acronyms and their expansions. They are rendered in \LaTeX{} with the \texttt{glossaries} package, which provides a rich set of commands for formatting and indexing.

\begin{code}{markdown}{}{}
\begin{Verbatim}[commandchars=\\\{\},breaklines, breakanywhere, commandchars=\\\{\}]
The HTML specification is maintained by the W3C.

*[HTML]: HyperText Markup Language
*[W3C]: World Wide Web Consortium
\end{Verbatim}
\end{code}
TeXSmith renders that snippet as:

\begin{code}{text}{}{}
\begin{Verbatim}[commandchars=\\\{\},breaklines, breakanywhere, commandchars=\\\{\}]
\PYZdl{} texsmith abbr.md
The \PYZbs{}acrshort\PYZob{}HTML\PYZcb{} specification is maintained by the \PYZbs{}acrshort\PYZob{}W3C\PYZcb{}.
\end{Verbatim}
\end{code}
Which displays as:

\begin{figure}[H]
    \centering

    \ifdefined\adjustbox

        \adjustbox{max width=\textwidth}{
            \includegraphics[width=\linewidth]{snippet-<HASH>.pdf}
        }

    \else

        \includegraphics[width=\linewidth]{snippet-<HASH>.pdf}

    \fi

\end{figure}
Of course, this also works on this \acrshort{HTML} site. Try hovering over the abbreviations.

Acronyms are collected automatically during the \LaTeX{} pass:

\begin{code}{text}{}{}
\begin{Verbatim}[commandchars=\\\{\},breaklines, breakanywhere, commandchars=\\\{\}]
\PYZdl{} uv run texsmith test.md \PYZhy{}tarticle 1\PYZgt{}/dev/null
\PYZdl{} rg newacronym build/test.tex
92:\PYZbs{}newacronym\PYZob{}HTML\PYZcb{}\PYZob{}HTML\PYZcb{}\PYZob{}HyperText Markup Language\PYZcb{}
93:\PYZbs{}newacronym\PYZob{}W3C\PYZcb{}\PYZob{}W3C\PYZcb{}\PYZob{}World Wide Web Consortium\PYZcb{}
\end{Verbatim}
\end{code}
\section{Front-matter glossary}\label{front-matter-glossary}

For longer documents you can declare acronyms in a structured \texttt{glossary:} section
in the YAML front matter. Each entry carries an explicit description and may be
attached to a group; TeXSmith renders one localised \texttt{\textbackslash{}printglossary} table per
group (in declaration order) followed by a default table for ungrouped entries.
The legacy \texttt{*[KEY]: …} body syntax keeps working and merges with the
front-matter entries.

\begin{code}{yaml}{}{}
\begin{Verbatim}[commandchars=\\\{\},breaklines, breakanywhere, commandchars=\\\{\}]
\PY{n+nn}{\PYZhy{}\PYZhy{}\PYZhy{}}
\PY{n+nt}{glossary}\PY{p}{:}
\PY{+w}{  }\PY{n+nt}{style}\PY{p}{:}\PY{+w}{ }\PY{l+lScalar+lScalarPlain}{long}\PY{+w}{           }\PY{c+c1}{\PYZsh{} default; any glossaries\PYZhy{}package style works}
\PY{+w}{  }\PY{n+nt}{groups}\PY{p}{:}
\PY{+w}{    }\PY{n+nt}{technique}\PY{p}{:}\PY{+w}{ }\PY{l+lScalar+lScalarPlain}{Acronymes}\PY{l+lScalar+lScalarPlain}{ }\PY{l+lScalar+lScalarPlain}{techniques}
\PY{+w}{    }\PY{n+nt}{institutionnel}\PY{p}{:}\PY{+w}{ }\PY{l+lScalar+lScalarPlain}{Acronymes}\PY{l+lScalar+lScalarPlain}{ }\PY{l+lScalar+lScalarPlain}{institutionnels}
\PY{+w}{  }\PY{n+nt}{entries}\PY{p}{:}
\PY{+w}{    }\PY{n+nt}{API}\PY{p}{:}
\PY{+w}{      }\PY{n+nt}{group}\PY{p}{:}\PY{+w}{ }\PY{l+lScalar+lScalarPlain}{technique}
\PY{+w}{      }\PY{n+nt}{description}\PY{p}{:}\PY{+w}{ }\PY{l+lScalar+lScalarPlain}{Application}\PY{l+lScalar+lScalarPlain}{ }\PY{l+lScalar+lScalarPlain}{Programming}\PY{l+lScalar+lScalarPlain}{ }\PY{l+lScalar+lScalarPlain}{Interface}
\PY{+w}{    }\PY{n+nt}{ONU}\PY{p}{:}
\PY{+w}{      }\PY{n+nt}{group}\PY{p}{:}\PY{+w}{ }\PY{l+lScalar+lScalarPlain}{institutionnel}
\PY{+w}{      }\PY{n+nt}{description}\PY{p}{:}\PY{+w}{ }\PY{l+lScalar+lScalarPlain}{Organisation}\PY{l+lScalar+lScalarPlain}{ }\PY{l+lScalar+lScalarPlain}{des}\PY{l+lScalar+lScalarPlain}{ }\PY{l+lScalar+lScalarPlain}{Nations}\PY{l+lScalar+lScalarPlain}{ }\PY{l+lScalar+lScalarPlain}{Unies}
\PY{+w}{    }\PY{n+nt}{DOI}\PY{p}{:}\PY{+w}{ }\PY{l+lScalar+lScalarPlain}{Digital}\PY{l+lScalar+lScalarPlain}{ }\PY{l+lScalar+lScalarPlain}{Object}\PY{l+lScalar+lScalarPlain}{ }\PY{l+lScalar+lScalarPlain}{Identifier}\PY{+w}{   }\PY{c+c1}{\PYZsh{} short form: ungrouped, description only}
\PY{n+nn}{\PYZhy{}\PYZhy{}\PYZhy{}}
\end{Verbatim}
\end{code}
The section is validated with pydantic, so unknown keys, missing descriptions,
or references to undeclared groups raise a clear error at conversion time. The
default acronym-table title follows the document language (it expands to
\texttt{\textbackslash{}acronymname} from the \texttt{glossaries} package, which is localised by \texttt{babel}).

\subsection{Automatic substitution and limitations}\label{automatic-substitution-and-limitations}

Unlike \LaTeX{}, TeXSmith does \textbf{not} require \texttt{\textbackslash{}gls\{…\}} / \texttt{\textbackslash{}Gls\{…\}} calls in the
source: the converter scans the body and replaces every \textbf{strict, case-sensitive}
match of an acronym key with \texttt{\textbackslash{}acrshort\{KEY\}}. As a consequence, casing helpers
such as \texttt{\textbackslash{}Gls}, \texttt{\textbackslash{}GLS}, \texttt{\textbackslash{}acrlong}, etc. are not synthesised --- the substitution
is the same regardless of where the acronym appears in the text. If you need
those forms, drop down to raw \LaTeX{} (e.g. with an explicit \texttt{\textbackslash{}Gls\{KEY\}} written
as a raw-\LaTeX{} inline).

\printglossary[type=\acronymtype]

\end{document}
